% Source: https://artofproblemsolving.com/community/c3h2222611_how_do_you_convert_fractions_into_different_bases_in_decimal_form
% Stable identity: aops-topic-2222611-q01
Multiply the fractional part repeatedly by the base; the integer parts are the digits.

\textbf{(a)} Base $2$:
\[\dfrac{1}{12}\cdot2=\dfrac{1}{6}\ (0),\quad\dfrac{1}{6}\cdot2=\dfrac{1}{3}\ (0),\quad\dfrac{1}{3}\cdot2=\dfrac{2}{3}\ (0),\quad\dfrac{2}{3}\cdot2=1+\dfrac{1}{3}\ (1),\quad\dfrac{1}{3}\cdot2=\dfrac{2}{3}\ (0),\ \ldots\]
From the third step on, the remainders alternate $\dfrac{2}{3},\dfrac{1}{3},\dfrac{2}{3},\ldots$, so the digits are $0,0,0,1,0,1,0,1,\ldots$:
\[\dfrac{1}{12}=\boxed{0.00\overline{01}_2}=0.000101010\ldots{}_2.\]
Check: $0.00\overline{01}_2=\dfrac{1}{4}\cdot0.\overline{01}_2=\dfrac{1}{4}\cdot\dfrac{1}{3}=\dfrac{1}{12}$, since $0.\overline{01}_2=\dfrac{1}{2^2-1}$.

\textbf{(b)} Base $3$:
\[\dfrac{1}{12}\cdot3=\dfrac{1}{4}\ (0),\quad\dfrac{1}{4}\cdot3=\dfrac{3}{4}\ (0),\quad\dfrac{3}{4}\cdot3=2+\dfrac{1}{4}\ (2),\quad\dfrac{1}{4}\cdot3=\dfrac{3}{4}\ (0),\ \ldots\]
The remainders alternate $\dfrac{3}{4},\dfrac{1}{4}$, so
\[\dfrac{1}{12}=\boxed{0.0\overline{02}_3}=0.00202\ldots{}_3.\]
Check: $0.0\overline{02}_3=\dfrac{1}{3}\cdot\dfrac{2}{3^2-1}=\dfrac{1}{3}\cdot\dfrac{2}{8}=\dfrac{1}{12}$.
